\documentclass[10pt,a4paper]{amsart}
\usepackage{amsmath,amssymb,amsthm,graphicx}
% Retained preamble for the Stacks Project chapter snapshot in this directory.
% The chapter's own preamble.tex is a copy of the shared Stacks Project preamble; the block
% below is sliced MECHANICALLY, line for line, from it. Nothing was added, removed or
% redefined; the wrapper class, the xr-hyper/hyperref wiring and the externaldocument list
% are omitted because the assembler supplies the document class.
%
% lines taken from preamble.tex: 3, 5, 9, 11, 12, 156, 157, 158, ... (51 lines total)
%
% package coverage restored from the chapter's own preamble (2026-10-13): xy, xy-2cell

\usepackage[all]{xy}
\xyoption{2cell}
\UseAllTwocells

\newenvironment{reference}{\comment}{\endcomment}
\newenvironment{slogan}{\comment}{\endcomment}
\newenvironment{history}{\comment}{\endcomment}
\theoremstyle{plain}
\newtheorem{theorem}[subsection]{Theorem}
\newtheorem{proposition}[subsection]{Proposition}
\newtheorem{lemma}[subsection]{Lemma}
\theoremstyle{definition}
\newtheorem{definition}[subsection]{Definition}
\newtheorem{example}[subsection]{Example}
\newtheorem{exercise}[subsection]{Exercise}
\newtheorem{situation}[subsection]{Situation}
\theoremstyle{remark}
\newtheorem{remark}[subsection]{Remark}
\newtheorem{remarks}[subsection]{Remarks}
\numberwithin{equation}{subsection}
\def\lim{\mathop{\mathrm{lim}}\nolimits}
\def\colim{\mathop{\mathrm{colim}}\nolimits}
\def\Spec{\mathop{\mathrm{Spec}}}
\def\Hom{\mathop{\mathrm{Hom}}\nolimits}
\def\Ext{\mathop{\mathrm{Ext}}\nolimits}
\def\SheafHom{\mathop{\mathcal{H}\!\mathit{om}}\nolimits}
\def\SheafExt{\mathop{\mathcal{E}\!\mathit{xt}}\nolimits}
\def\Sch{\mathit{Sch}}
\def\Mor{\mathop{\mathrm{Mor}}\nolimits}
\def\Ob{\mathop{\mathrm{Ob}}\nolimits}
\def\Sh{\mathop{\mathit{Sh}}\nolimits}
\def\NL{\mathop{N\!L}\nolimits}
\def\CH{\mathop{\mathrm{CH}}\nolimits}
\def\proetale{{pro\text{-}\acute{e}tale}}
\def\etale{{\acute{e}tale}}
\def\QCoh{\mathit{QCoh}}
\def\Ker{\mathop{\mathrm{Ker}}}
\def\Im{\mathop{\mathrm{Im}}}
\def\Coker{\mathop{\mathrm{Coker}}}
\def\Coim{\mathop{\mathrm{Coim}}}
\DeclareMathSymbol{\boxtimes}{\mathbin}{AMSa}{"02}
\def\QCohstack{\mathcal{QC}\!\mathit{oh}}
\def\Cohstack{\mathcal{C}\!\mathit{oh}}
\def\Spacesstack{\mathcal{S}\!\mathit{paces}}
\def\Quotfunctor{\mathrm{Quot}}
\def\Hilbfunctor{\mathrm{Hilb}}
\def\Curvesstack{\mathcal{C}\!\mathit{urves}}
\def\Polarizedstack{\mathcal{P}\!\mathit{olarized}}
\def\Complexesstack{\mathcal{C}\!\mathit{omplexes}}
\def\Pic{\mathop{\mathrm{Pic}}\nolimits}
\def\Picardstack{\mathcal{P}\!\mathit{ic}}
\def\Picardfunctor{\mathrm{Pic}}
\def\Deformationcategory{\mathcal{D}\!\mathit{ef}}
\usepackage{cleveref}
\crefname{originallemma}{original lemma}{original lemmas}
\Crefname{originallemma}{Original lemma}{Original lemmas}
\makeatletter
\newlabel{descent-lemma-open-fpqc-covering}{{03N0}{source}{}{sourcepointer-1}{}}
\newlabel{descent-lemma-open-fpqc-covering@cref}{{[originallemma][0][]03N0}{[1][0][]source}{}{}{}}
\makeatother
\crefname{originallemma}{original lemma}{original lemmas}
\Crefname{originallemma}{Original lemma}{Original lemmas}
\makeatletter
\newlabel{groupoids-lemma-constructing-invariant-opens}{{03LO}{source}{}{sourcepointer-2}{}}
\newlabel{groupoids-lemma-constructing-invariant-opens@cref}{{[originallemma][0][]03LO}{[1][0][]source}{}{}{}}
\makeatother
\crefname{originallemma}{original lemma}{original lemmas}
\Crefname{originallemma}{Original lemma}{Original lemmas}
\makeatletter
\newlabel{groupoids-lemma-quotient-pre-equivalence-relation-restrict}{{045Z}{source}{}{sourcepointer-3}{}}
\newlabel{groupoids-lemma-quotient-pre-equivalence-relation-restrict@cref}{{[originallemma][0][]045Z}{[1][0][]source}{}{}{}}
\makeatother
\crefname{originallemma}{original lemma}{original lemmas}
\Crefname{originallemma}{Original lemma}{Original lemmas}
\makeatletter
\newlabel{groupoids-lemma-restrict-relation}{{02V8}{source}{}{sourcepointer-4}{}}
\newlabel{groupoids-lemma-restrict-relation@cref}{{[originallemma][0][]02V8}{[1][0][]source}{}{}{}}
\makeatother
\crefname{originallemma}{original lemma}{original lemmas}
\Crefname{originallemma}{Original lemma}{Original lemmas}
\makeatletter
\newlabel{morphisms-lemma-fppf-open}{{01UA}{source}{}{sourcepointer-5}{}}
\newlabel{morphisms-lemma-fppf-open@cref}{{[originallemma][0][]01UA}{[1][0][]source}{}{}{}}
\makeatother
\title{A restricted flat, locally finitely presented pre-relation has open quotient map}
\author{The Stacks Project Authors}
\date{stacks-project:spaces.tex:local-snapshot-20261012; GNU FDL 1.2 or later}
\begin{document}
\maketitle
\setcounter{subsection}{1}
\noindent\textit{Editorial note.} All delivered mathematical text is the author's own, copied from the pinned Stacks Project chapter file spaces.tex; no statement or proof step was rewritten, reordered or removed.\par\medskip
\noindent\textit{Editorial note.} The proof contains two xymatrix displays; they are part of the counted proof span and are preserved verbatim, with the chapter's own Xy-pic setup retained in the preamble.\par\medskip
\noindent\textit{Editorial note.} The five references of the counted spans are recorded as non-load-bearing navigation pointers; the argument that consumes them stays inside the counted proof span.\par\medskip
\noindent\textit{Editorial note.} Preamble: sources/stacks-spaces/retained-preamble.tex, sliced mechanically from the chapter's own preamble.tex.\par\medskip
\noindent\textit{Allowed standard prerequisites (editorial).}\begin{itemize}
\item Restriction of an equivalence relation along a morphism is an equivalence relation. Groupoids, Lemma groupoids-lemma-restrict-relation, used at proof lines 1544-1546. Gives the equivalence-relation hypothesis for R'.
\item Flat and locally finitely presented morphisms are universally open. Morphisms, Lemma morphisms-lemma-fppf-open, used at proof lines 1546-1549. Makes W\textasciicircum{}1 and W open in U.
\item The invariant open attached to an equivalence relation. Groupoids, Lemma groupoids-lemma-constructing-invariant-opens, used at proof lines 1550-1554. Supplies t(s\textasciicircum{}\{-1\}(W)) = W.
\item Injectivity of U'/R' -> U/R for the restricted relation. Groupoids, Lemma groupoids-lemma-quotient-pre-equivalence-relation-restrict, used at proof lines 1557-1560. Turns the claim into representing the fibre product by an open subscheme.
\item Openness descends through an fpqc covering. Descent, Lemma descent-lemma-open-fpqc-covering, used at proof lines 1594-1597. Produces the open W\_T subset T with phi\_j\textasciicircum{}\{-1\}(W\_T) = W\_j.
\end{itemize}\medskip
\ifdefined\hypertarget\hypertarget{sourcepointer-1}{}\fi
\noindent\textit{Source reference.} Original lemma 03N0: focused-seeds/AGP/sources/stacks-snapshot-20261011/tags.txt, source lines 4768--4768. Navigation pointer into the Stacks Project; the standard fact it supplies is applied inside the counted proof span and the argument resting on it stays there.\par\medskip
\ifdefined\hypertarget\hypertarget{sourcepointer-2}{}\fi
\noindent\textit{Source reference.} Original lemma 03LO: focused-seeds/AGP/sources/stacks-snapshot-20261011/tags.txt, source lines 4720--4720. Navigation pointer into the Stacks Project; the standard fact it supplies is applied inside the counted proof span and the argument resting on it stays there.\par\medskip
\ifdefined\hypertarget\hypertarget{sourcepointer-3}{}\fi
\noindent\textit{Source reference.} Original lemma 045Z: focused-seeds/AGP/sources/stacks-snapshot-20261011/tags.txt, source lines 5454--5454. Navigation pointer into the Stacks Project; the standard fact it supplies is applied inside the counted proof span and the argument resting on it stays there.\par\medskip
\ifdefined\hypertarget\hypertarget{sourcepointer-4}{}\fi
\noindent\textit{Source reference.} Original lemma 02V8: focused-seeds/AGP/sources/stacks-snapshot-20261011/tags.txt, source lines 3759--3759. Navigation pointer into the Stacks Project; the standard fact it supplies is applied inside the counted proof span and the argument resting on it stays there.\par\medskip
\ifdefined\hypertarget\hypertarget{sourcepointer-5}{}\fi
\noindent\textit{Source reference.} Original lemma 01UA: focused-seeds/AGP/sources/stacks-snapshot-20261011/tags.txt, source lines 2403--2403. Navigation pointer into the Stacks Project; the standard fact it supplies is applied inside the counted proof span and the argument resting on it stays there.\par\medskip
\begin{lemma}
\label{lemma-finding-opens}
Let $S$ be a scheme.
Let $U$ be a scheme over $S$.
Let $j = (s, t) : R \to U \times_S U$ be a pre-relation.
Let $g : U' \to U$ be a morphism.
Assume
\begin{enumerate}
\item $j$ is an equivalence relation,
\item $s, t : R \to U$ are surjective, flat and
locally of finite presentation,
\item $g$ is flat and locally of finite presentation.
\end{enumerate}
Let $R' = R|_{U'}$ be the restriction of $R$ to $U'$. Then
$U'/R' \to U/R$ is representable, and is an open immersion.
\end{lemma}

\begin{proof}
By Groupoids, Lemma \ref{groupoids-lemma-restrict-relation}
the morphism $j' = (s', t') : R' \to U' \times_S U'$
defines an equivalence relation. Since $g$ is flat and locally of
finite presentation we see that $g$ is universally open as well
(Morphisms, Lemma \ref{morphisms-lemma-fppf-open}).
For the same reason $s, t$ are universally open as well.
Let $W^1 = g(U') \subset U$, and let $W = t(s^{-1}(W^1))$.
Then $W^1$ and $W$ are open in $U$. Moreover, as $j$ is an
equivalence relation we have $t(s^{-1}(W)) = W$ (see
Groupoids, Lemma \ref{groupoids-lemma-constructing-invariant-opens}
for example).

\medskip\noindent
By
Groupoids,
Lemma \ref{groupoids-lemma-quotient-pre-equivalence-relation-restrict}
the map of sheaves $F' = U'/R' \to F = U/R$ is injective.
Let $a : T \to F$ be a morphism from a scheme into $U/R$.
We have to show that $T \times_F F'$ is representable
by an open subscheme of $T$.

\medskip\noindent
The morphism $a$ is given by the following data:
an fppf covering $\{\varphi_j : T_j \to T\}_{j \in J}$ of $T$ and
morphisms $a_j : T_j \to U$ such that the maps
$$
a_j \times a_{j'} :
T_j \times_T T_{j'}
\longrightarrow
U \times_S U
$$
factor through $j : R \to U \times_S U$ via some (unique) maps
$r_{jj'} : T_j \times_T T_{j'} \to R$. The system
$(a_j)$ corresponds to $a$ in the sense that the diagrams
$$
\xymatrix{
T_j \ar[r]_{a_j} \ar[d] & U \ar[d] \\
T \ar[r]^a & F
}
$$
commute.

\medskip\noindent
Consider the open subsets $W_j = a_j^{-1}(W) \subset T_j$.
Since $t(s^{-1}(W)) = W$ we see that
$$
W_j \times_T T_{j'} =
r_{jj'}^{-1}(t^{-1}(W)) = r_{jj'}^{-1}(s^{-1}(W)) =
T_j \times_T W_{j'}.
$$
By
Descent, Lemma \ref{descent-lemma-open-fpqc-covering}
this means there exists an open
$W_T \subset T$ such that $\varphi_j^{-1}(W_T) = W_j$ for all $j \in J$.
We claim that $W_T \to T$ represents $T \times_F F' \to T$.

\medskip\noindent
First, let us show that $W_T \to T \to F$ is an element of
$F'(W_T)$. Since $\{W_j \to W_T\}_{j \in J}$ is an
fppf covering of $W_T$, it is enough to show that
each $W_j \to U \to F$ is an element of $F'(W_j)$ (as $F'$ is a sheaf
for the fppf topology). Consider the commutative diagram
$$
\xymatrix{
W'_j \ar[rr] \ar[dd] \ar[rd] & & U' \ar[d]^g \\
& s^{-1}(W^1) \ar[r]_s \ar[d]^t & W^1 \ar[d] \\
W_j \ar[r]^{a_j|_{W_j}} & W \ar[r] & F
}
$$
where $W'_j = W_j \times_W s^{-1}(W^1) \times_{W^1} U'$.
Since $t$ and $g$ are surjective, flat and locally of finite
presentation, so is $W'_j \to W_j$. Hence the restriction of
the element $W_j \to U \to F$ to $W'_j$ is an element of $F'$
as desired.

\medskip\noindent
Suppose that $f : T' \to T$ is a morphism of schemes
such that $a|_{T'} \in F'(T')$. We have to show that
$f$ factors through the open $W_T$. Since
$\{T' \times_T T_j \to T'\}$ is an fppf covering of $T'$
it is enough to show each $T' \times_T T_j \to T$
factors through $W_T$. Hence we may assume $f$ factors
as $\varphi_j \circ f_j : T' \to T_j \to T$ for some $j$.
In this case the condition $a|_{T'} \in F'(T')$ means that there exists
some fppf covering $\{\psi_i : T'_i \to T'\}_{i \in I}$ and some
morphisms $b_i : T'_i \to U'$ such that
$$
\xymatrix{
T'_i \ar[r]_{b_i} \ar[d]_{f_j \circ \psi_i} & U' \ar[r]_g & U \ar[d] \\
T_j \ar[r]^{a_j} & U \ar[r] & F
}
$$
is commutative. This commutativity means that there exists a
morphism $r'_i : T'_i \to R$ such that
$t \circ r'_i = a_j \circ f_j \circ \psi_i$, and
$s \circ r'_i = g \circ b_i$. This implies that
$\Im(f_j \circ \psi_i) \subset W_j$ and we win.
\end{proof}

\end{document}
